\documentclass[10pt,a4paper]{article}
\usepackage{amsmath,amssymb,amsthm,graphicx}
\PassOptionsToPackage{dvipsnames}{xcolor}
% !TeX root = print.tex
% Put any macro and import needed for the project here.
% This will be used by both the web and print versions of the blueprint.
% This file is not meant to be built. Build src/web.tex or src/print.text instead.

\usepackage{mathtools}
\usepackage{amssymb}
\usepackage{amsthm}
\usepackage{amsmath}
\usepackage{graphicx}
\usepackage{xcolor}
\usepackage{enumerate}
% \usepackage{showlabels} % not compatible with Plastex
\usepackage[normalem]{ulem}

\newcommand{\rs}[1]{{\color{blue}  RS: #1.}}
\newcommand{\lars}[1]{{\color{red}  LB: #1.}}
\newcommand{\asgar}[1]{{\color{green} AJ: #1}} %\color{TealBlue}
\newcommand{\ct}[1]{{\color{purple}  CT: #1}}


\theoremstyle{plain}
\newtheorem{theorem}{Theorem}[subsection]
\newtheorem{lemma}[theorem]{Lemma}
\newtheorem{proposition}[theorem]{Proposition}
\newtheorem{cor}[theorem]{Corollary}
\theoremstyle{definition}
\newtheorem{definition}[theorem]{Definition}
\newtheorem{remark}[theorem]{Remark}
\newtheorem{example}[theorem]{Example}
\newtheorem{examples}[theorem]{Examples}
% \numberwithin{section}{chapter}
%\numberwithin{subsection}{section}
%\numberwithin{subsubsection}{subsection}
\numberwithin{equation}{subsection}

\usepackage[hidelinks, colorlinks=false, hypertexnames=false]{hyperref}


\newcommand{\N}{{\mathbb{N}}}
\newcommand{\Z}{{\mathbb{Z}}}
\newcommand{\Q}{{\mathbb{Q}}}
\newcommand{\R}{{\mathbb{R}}}

\DeclareMathOperator{\ch}{\operatorname{ch}}
\DeclareMathOperator{\dens}{\operatorname{dens}}
\DeclareMathOperator{\supp}{\operatorname{supp}}
\DeclareMathOperator{\tp}{\operatorname{top}}
\DeclareMathOperator{\im}{\operatorname{im}}
\DeclareMathOperator{\Lip}{\operatorname{Lip}}
\DeclareMathOperator{\bd}{\operatorname{bd}}
\DeclareMathOperator*{\esssup}{\operatorname{ess\,sup}}

\newcommand{\fp}{{\mathfrak p}}
\newcommand{\fP}{{\mathfrak P}}
\newcommand{\fu}{{\mathfrak u}}
\newcommand{\fU}{{\mathfrak U}}
\newcommand{\fv}{{\mathfrak v}}
\newcommand{\fq}{{\mathfrak q}}
\newcommand{\fQ}{{\mathfrak Q}}
\newcommand{\fT}{{\mathfrak T}}
\newcommand{\fL}{{\mathfrak L}}
\newcommand{\fC}{{\mathfrak C}}
\newcommand{\pc}{{\mathrm{c}}}
\newcommand{\ps}{{\mathrm{s}}}
\newcommand{\AD}{{\bf s}}
\newcommand{\fc}{{\Omega}}
\newcommand{\borel}{{\mathcal{E}}}
\newcommand{\borelb}{{\mathcal{G}}}
\newcommand{\scI}{{\mathcal{I}}} % note: renamed because it conflicted with another command
\newcommand{\tQ}{{Q}}
\newcommand{\mfa}{{\vartheta}}
\newcommand{\mfb}{{\theta}}
\newcommand{\Mf}{{\Theta}}
\newcommand{\fcc}{{\mathcal{Q}}}
\usepackage[a4paper,landscape,margin=18mm]{geometry}
\usepackage{cleveref}
\newcommand{\lean}[1]{}
\newcommand{\proves}[1]{}
\newcommand{\uses}[1]{}
\newcommand{\leanok}{}
\def\C{\mathbb{C}}
\usepackage{xurl}
\usepackage{unicode-math}
\setmathfont{latinmodern-math.otf}
\makeatletter
\renewcommand{\author}[1]{\gdef\@author{\parbox{0.94\textwidth}{\centering #1}}}
\makeatother
\title{Almost-everywhere Fourier convergence for continuous periodic functions}
\author{Lars Becker, María Inés de Frutos-Fernández, Leo Diedering, Floris van Doorn, Sébastien Gouëzel, Asgar Jamneshan, Evgenia Karunus, Edward van de Meent, Pietro Monticone, Jasper Mulder-Sohn, Jim Portegies, Joris Roos, Michael Rothgang, Rajula Srivastava, James Sundstrom, Jeremy Tan, Christoph Thiele}
\date{github:fpvandoorn-carleson:12fb202e6dba84aa57a9d39c7d019dc344902912; Apache-2.0}
\begin{document}
\maketitle
\noindent\textit{Editorial note.} Complete author mathematical source spans are preserved without generated proof text. Blueprint Lean/dependency display metadata is suppressed; theorem numbering is local.\par\medskip
\noindent\textit{Editorial note.} The original project paper discloses minor AI use for retrieval, typo suggestions, and one later Aristotle-assisted code PR; source disclosures are retained. No mathematical correctness audit is claimed.\par\medskip
\noindent\textit{Editorial note.} Formatting-only author block wraps the complete original17-author list within the page.\par\medskip
\noindent\textit{Allowed standard prerequisites (editorial).}\begin{itemize}
\item Standard measure theory, dominated/monotone convergence, Fourier Hilbert-space identities and elementary doubling-space facts are cited as in the author Blueprint.
\item For the classical theorem the metric-space Carleson theorem is an explicitly stated named prerequisite, with its separate complete proof retained in the companion material.
\end{itemize}\medskip
\subsection{Proof of Cotlar's Inequality}
\label{subsec-cotlar}


\section{Proof of Metric Space Carleson, overview}
\label{overviewsection}


With Euler's formula to represent pure frequencies in mind, a trigonometric polynomial can be expressed as
\begin{equation}\label{eq:trig-series}
    S_N(x):= \sum_{n=-N}^N c_n e^{inx} \ .
\end{equation}
The Fourier series is then defined as the limit $f$ of such a sequence $S_N$
as $N$ tends to $\infty$. Fourier's vision to represent rather general
functions raises two fundamental questions. The first question is to identify
the appropriate choice of coefficients $c_n$ to use to represent a given $f$.
The second question addresses the convergence of $S_N$ to $f$.

The first question has a fairly canonical and standard answer, provided by the Fourier integral formula:
\begin{equation}\label{eq:fourier-coefficients}
    c_n:=\widehat{f}_n:=\frac 1{2\pi}\int_0^{2\pi}f(x) e^{- i nx}\, dx,
\end{equation}

\subsection{Statement of the metric space Carleson theorems}

Let
\begin{equation}
    a\ge 4
\end{equation}
be a natural number that as it gets larger will allow for more general applications of Theorem \ref{metric-space-Carleson}
but will worsen the constants in the estimates.


A doubling metric measure space $(X,\rho,\mu, a)$ is a complete
and locally compact metric space $(X,\rho)$
equipped with a non-zero locally finite Borel measure $\mu$ that satisfies the doubling condition that for all $x\in X$ and all $R>0$ we have
\begin{equation}\label{doublingx}
    \mu(B(x,2R))\le 2^a\mu(B(x,R))\,,
\end{equation}
where we have denoted by $B(x,R)$ the open ball of radius $R$ centered at $x$:
\begin{equation}\label{eq-define-ball}
 B(x,R):=\{y\in X: \rho(x,y)<R\}. \end{equation}
In a doubling metric measure space the closed balls are compact and $\mu$ is positive on all non-empty open sets.

A collection $\Mf$ of real valued continuous functions on the doubling metric measure space $(X,\rho,\mu,a)$ is called compatible,
if there is a point $o\in X$ where all the functions are equal to $0$,
and if there exists for each ball $B \subset X$ a metric $d_B$ on $\Mf$,
such that the following five properties \eqref{osccontrol}, \eqref{firstdb}, \eqref{monotonedb}, \eqref{seconddb}, and \eqref{thirddb} are satisfied.
For every ball $B \subset X$
\begin{equation}\label{osccontrol}
    \sup_{x,y\in B}|\mfa(x)-{\mfa(y)}- \mfb(x)+{\mfb(y)}| \le d_{B}(\mfa,\mfb)\,.
\end{equation}
For any two balls $B_1=B(x_1,R)$, $B_2= B(x_2,2R)$ in $X$ with $x_1\in B_2$ and any $\mfa,\mfb\in \Mf$,
\begin{equation}\label{firstdb}
    d_{B_2}(\mfa,\mfb)\le 2^a d_{B_1}(\mfa,\mfb) .
\end{equation}
For any two balls $B_1, B_2$ in $X$ with $B_1 \subset B_2$ and any $\mfa, \mfb \in \Mf$
\begin{equation}\label{monotonedb}
    d_{B_1}(\mfa,\mfb) \le d_{B_2}(\mfa, \mfb)
\end{equation}
and for any two balls
$B_1=B(x_1,R)$, $B_2= B(x_2,2^aR)$
with $B_1\subset B_2$, and $\mfa,\mfb\in \Mf$,
\begin{equation}\label{seconddb}
    2d_{B_1}(\mfa,\mfb)
\le d_{B_2}(\mfa,\mfb) .
\end{equation}
For every ball $B$ in $X$ and every $d_B$-ball $\tilde B$ of radius $2R$ in $\Mf$, there is a collection $\mathcal{B}$ of
at most $2^a$ many $d_B$-balls of radius $R$ covering $\tilde B$, that is,
\begin{equation}\label{thirddb}
    \tilde B\subset \bigcup \mathcal{B}.
\end{equation}


Further, a compatible collection $\Mf$ is called cancellative, if
for any ball $B$ in $X$ of radius $R$, any Lipschitz function $\varphi: X\to \C$
supported on $B$, and any $\mfa,\mfb\in \Mf$ we have
\begin{equation}
    \label{eq-vdc-cond}
    |\int_B e(\mfa(x)-{\mfb(x)}) \varphi(x) d\mu(x)|\le 2^a \mu(B)\|\varphi\|_{\Lip(B)}
(1+d_B(\mfa,\mfb))^{-\frac{1}{a}},
\end{equation}
where $\|\cdot\|_{\Lip(B)}$ denotes the inhomogeneous Lipschitz norm on $B$:
$$
    \|\varphi\|_{\Lip(B)} = \sup_{x \in B} |\varphi(x)| + R \sup_{x,y \in B, x \neq y} \frac{|\varphi(x) - \varphi(y)|}{\rho(x,y)}\,.
$$

A one-sided Calder\'on--Zygmund kernel $K$ on the doubling metric measure space $(X, \rho, \mu, a)$ is a measurable function
\begin{equation}\label{eqkernel0}
  K:X\times X\to \mathbb{C}
\end{equation}
such that for all $x,y',y\in X$ with $x\neq y$, we have
\begin{equation}\label{eqkernel-size}
    |K(x,y)| \leq \frac{2^{a^3}}{V(x,y)}
\end{equation}
and if $2\rho(y,y') \leq \rho(x,y)$, then
\begin{equation}
  \label{eqkernel-y-smooth}
  |K(x,y) - K(x,y')| \leq \left(\frac{\rho(y,y')}{\rho(x,y)}\right)^{\frac{1}{a}}\frac{2^{a^3}}{V(x,y)},
\end{equation}
where \[V(x,y):=\mu(B(x,\rho(x,y))).\]
Define the maximally truncated non-tangential singular integral $T_{*}$ associated with $K$ by
\begin{equation}
    \label{def-tang-unm-op}
    T_{*}f(x):=\sup_{R_1 < R_2} \sup_{\rho(x,x')<R_1} \left|\int_{R_1< \rho(x',y) < R_2} K(x',y) f(y) \, \mathrm{d}\mu(y) \right|\,.
\end{equation}
We define the generalized Carleson operator $T$ by
\begin{equation}
    \label{def-main-op}
    Tf(x):=\sup_{\mfa\in\Mf} \sup_{0 < R_1 < R_2}\left| \int_{R_1 < \rho(x,y) < R_2} K(x,y) f(y) e(\mfa(y)) \, \mathrm{d}\mu(y) \right|\, ,
\end{equation}
where $e(r)=e^{ir}$.

Our first main result is the following restricted weak type estimate for $T$ in the range $1<q\le 2$,
which by interpolation techniques recovers $L^q$ estimates for the open range $1<q<2$.
\begin{theorem}[metric space Carleson]
\label{metric-space-Carleson}
\uses{linearised-metric-Carleson,int-continuous}
\leanok
\lean{metric_carleson}
    For all integers $a \ge 4$ and real numbers $1<q\le 2$ the following holds.
    Let $(X,\rho,\mu,a)$ be a doubling metric measure space.
    Let $\Mf$ be a cancellative compatible collection of functions and let $K$ be a one-sided Calder\'on--Zygmund kernel on $(X,\rho,\mu,a)$.
    Assume that for every bounded measurable function $g$ on $X$ supported on a set of finite measure we have
    \begin{equation}\label{nontanbound}
        \|T_{*}g\|_{2} \leq 2^{a^3} \|g\|_2\,,
    \end{equation}
    where $T_{*}$ is defined in \eqref{def-tang-unm-op}.
    Then for all Borel sets $F$ and $G$ in $X$ and all Borel functions $f:X\to \C$ with
    $|f|\le \mathbf{1}_F$, we have, with $T$ defined in \eqref{def-main-op},
    \begin{equation}
        \label{resweak}
        \left|\int_{G} T f \, \mathrm{d}\mu\right| \leq \frac{2^{443a^3}}{(q-1)^6} \mu(G)^{1-\frac{1}{q}} \mu(F)^{\frac{1}{q}}\, .
    \end{equation}
\end{theorem}

In some applications, such as the Walsh-case of Carleson's theorem, the kernel $K$ naturally depends also on the modulation functions $\mfa$.
The fact that we don't assume H\"older-continuity of the kernel $K$ in the first argument allows us to also capture this situation.

We now state another Theorem with a slightly weaker replacement of the assumption \eqref{nontanbound}.
It implies the Walsh-case of Carleson's theorem.
For a Borel function $\tQ:X\to \Mf$, and $\mfa \in \Mf$ and $x\in X$ define
\begin{equation}
    R_{\tQ}(\mfa,x)=\sup\{r>0:d_{B(x,r)}(\mfa, \tQ(x))<1\}
\end{equation}
and define further
\begin{equation}
    \label{def-lin-star-op}
    T_{\tQ}^\mfa f(x):=\sup_{R_1<R_2} \ \sup_{\rho(x,x')<R_1}
    \left|\int_{R_1< \rho(x',y) < \min\{R_2, R_{\tQ}(\mfa,x')\}} K(x',y) f(y) \, \mathrm{d}\mu(y) \right|
\end{equation}
Define further the linearized generalized Carleson operator $T_\tQ$ by
\begin{equation}
    \label{def-lin-main-op}
    T_\tQ f(x):= \sup_{0 < R_1 < R_2}\left| \int_{R_1 < \rho(x,y) < R_2} K(x,y) f(y) e(\tQ(x)(y)) \, \mathrm{d}\mu(y) \right|\, ,
\end{equation}
where again $e(r)=e^{ir}$.

\begin{theorem}[linearised metric Carleson]
\label{linearised-metric-Carleson}
\uses{R-truncation,int-continuous}
\leanok
\lean{linearized_metric_carleson}
    For all integers $a \ge 4$ and real numbers $1<q\le 2$ the following holds.
    Let $(X,\rho,\mu,a)$ be a doubling metric measure space. Let $\Mf$ be a
    cancellative compatible collection of functions.
    Let $\tQ:X\to \Mf$ be a Borel function with finite range.
    Let $K$ be a one-sided Calder\'on--Zygmund kernel on $(X,\rho,\mu,a)$. Assume that for every $\mfa\in \Mf$ and every bounded measurable function $g$ on $X$ supported on a set of finite measure we have
    \begin{equation}\label{linnontanbound}
        \|T_{\tQ}^\mfa g\|_{2} \leq 2^{a^3} \|g\|_2\,,
    \end{equation}
    where $T_{\tQ}^\mfa$ is defined in \eqref{def-lin-star-op}.
    Then for all Borel sets $F$ and $G$ in $X$ and all Borel functions $f:X\to \C$ with
    $|f|\le \mathbf{1}_F$, we have, with $T_\tQ$ defined in \eqref{def-lin-main-op},
    \begin{equation}
        \label{linresweak}
        \left|\int_{G} T_\tQ f \, \mathrm{d}\mu\right| \le \frac{2^{443a^3}}{(q-1)^6} \mu(G)^{1-\frac{1}{q}} \mu(F)^{\frac{1}{q}}\, .
    \end{equation}
\end{theorem}

\section{Two-sided Metric Space Carleson}

We prove a variant of \Cref{metric-space-Carleson} for a two-sided Calder\'on--Zygmund kernel on the doubling metric measure space $(X,\rho,\mu,a)$, i.e. a one-sided Calder\'on--Zygmund kernel $K$ which additionally satisfies for all $x,x',y\in X$ with $x\neq y$ and $2\rho(x,x') \leq \rho(x,y)$,
\begin{equation}
    \label{eqkernel-x-smooth}
    |K(x,y) - K(x',y)| \leq \left(\frac{\rho(x,x')}{\rho(x,y)}\right)^{\frac{1}{a}}\frac{2^{a^3}}{V(x,y)}.
\end{equation}
%fixme: Should it be V(y,x) rather?
By the additional regularity, we can weaken the assumption \eqref{nontanbound} to a family of operators that is easier to work with in applications.
Namely, for $r > 0$, $x\in X$, and a bounded, measurable function $f:X\to\C$ supported on a set of finite measure, we define
\begin{equation}
\label{def-T-r}
T_r f(x):= \int_{r\le\rho(x,y)} K(x,y) f(y) \, d\mu(y) = \int_{X\setminus B(x,r)} K(x,y) f(y) \, d\mu(y).
\end{equation}
\begin{theorem}[two-sided metric space Carleson]
    \label{two-sided-metric-space-Carleson}
    \uses{metric-space-Carleson, nontangential-from-simple}
    \leanok
    \lean{two_sided_metric_carleson}
        For all  integers $a \ge  4$ and real numbers $1<q\le 2$
        the following holds.
        Let $(X,\rho,\mu,a)$ be a doubling metric measure space. Let  $\Mf$ be a
        cancellative compatible  collection of functions and let $K$ be a two-sided Calder\'on--Zygmund kernel on $(X,\rho,\mu,a)$. Assume  that for every bounded measurable function $g$ on $X$ supported on a set of finite measure and all $r>0$ we have
      \begin{equation}\label{two-sided-Hr-bound-assumption}
            \|T_r g\|_{2} \leq 2^{a^3} \|g\|_2\,.
        \end{equation}
        Then for all Borel sets $F$ and $G$ in $X$ and
        all Borel functions $f:X\to \C$ with
        $|f|\le \mathbf{1}_F$, we have, with $T$ defined in \eqref{def-main-op},
      \begin{equation}
        \label{two-sided-resweak}
            \left|\int_{G} T f \, \mathrm{d}\mu\right| \leq \frac{2^{474a^3}}{(q-1)^6} \mu(G)^{1-\frac{1}{q}} \mu(F)^{\frac{1}{q}}\, .
        \end{equation}
\end{theorem}

For the remainder of this chapter, fix an integer $a\ge 4$, a doubling metric measure space $(X,\rho,\mu,a)$ and a two-sided Calder\'on--Zygmund kernel $K$ as in \Cref{two-sided-metric-space-Carleson}.

The following lemma is proved in \Cref{subsec-cotlar}.
\begin{lemma}[nontangential-from-simple]\label{nontangential-from-simple}
    \leanok
    \lean{nontangential_from_simple}
    \uses{simple-nontangential-operator, nontangential-operator-boundary}
    Assume \eqref{two-sided-Hr-bound-assumption} holds.
    Then, for every bounded measurable function $g : X \to \C$ supported on a set of finite measure we have
    \begin{equation}\label{concretetstarbound}
        \|T_*g\|_2\le 2^{3a^3}\|g\|_2.
    \end{equation}
\end{lemma}

\section{Proof of The Classical Carleson Theorem}

The convergence of partial Fourier sums is proved in
\Cref{10classical} in two steps. In the first step,
we establish convergence on a suitable dense subclass of functions. We choose smooth functions as subclass, the convergence is stated in \Cref{convergence-for-smooth} and proved in \Cref{10smooth}.
In the second step, one controls the relevant error of approximating a general function by a function in
the subclass. This is stated in \Cref{control-approximation-effect} and proved
in \Cref{10difference}.
The proof relies on a bound on the real Carleson maximal operator stated in \Cref{real-Carleson} and proved in \Cref{10carleson}, which involves showing that the real line fits into the setting of \Cref{overviewsection}.
This latter proof refers to the two-sided variant of the Carleson \Cref{two-sided-metric-space-Carleson}. Two assumptions in \Cref{metric-space-Carleson} require more work. The boundedness of the operator $T_r$ defined in \eqref{def-T-r} is established in \ref{Hilbert-strong-2-2}. This lemma is proved in \Cref{10hilbert}.
The cancellative property is verified by \Cref{van-der-Corput}, which is proved in \Cref{10vandercorput}.
Several further auxiliary lemmas are stated and proved in \Cref{10classical}, the proof of one of these auxiliary lemmas, \Cref{spectral-projection-bound}, is done in \Cref{10projection}.

All subsections past \Cref{10classical} are mutually independent.



\subsection{The classical Carleson theorem}
\label{10classical}

Let a uniformly continuous $2\pi$-periodic function $f:\R\to \mathbb{C}$ and $\epsilon>0$ be given.
Let
\begin{equation}
    C_{a,q} := \frac{2^{474a^3}}{(q-1)^6}
\end{equation}
denote the constant from \Cref{two-sided-metric-space-Carleson}.
Define
\begin{equation}
    \epsilon' := \frac {\epsilon} {4 C_\epsilon} \,,
\end{equation}
where
%fixme: replace by C(\epsilon) ?
\begin{equation*}
    C_\epsilon = \left(\frac{8}{\pi\epsilon}\right)^\frac{1}{2} C_{4,2} + \pi \,.
\end{equation*}
Since $f$ is continuous and periodic, $f$ is uniformly continuous.
Thus, there is a $0<\delta<\pi$
such that for all $x,x' \in \R$ with $|x-x'|\le \delta$
we have
\begin{equation}\label{uniconbound}
|f(x)-f(x')|\le \epsilon' \, .
\end{equation}
Define
\begin{equation}\label{def-fzero}
f_0:=f \ast \phi_\delta,
\end{equation}
where $\phi_\delta$ is a nonnegative smooth bump function with $\supp (\phi_\delta) \subset (-\delta, \delta)$ and $\int _\R \phi_\delta (x) \, dx = 1$.

\begin{lemma}[smooth approximation]
\label{smooth-approximation}
\leanok
\lean{close_smooth_approx_periodic}
    The function $f_0$ is $2\pi$-periodic.
    The function $f_0$ is smooth (and therefore measurable).
    The function $f_0$ satisfies for all $x\in \R$:
    \begin{equation}\label{eq-ffzero}
    |f(x)-f_0(x)|\le \epsilon' \, ,
    \end{equation}
\end{lemma}

\begin{proof}
    \leanok
    Periodicity follows directly from the definitions. The other properties are part of the Lean library.
\end{proof}

We prove in \Cref{10smooth}:
\begin{lemma}[convergence for smooth]
\label{convergence-for-smooth}
\uses{convergence-for-twice-contdiff}
\leanok
\lean{fourierConv_ofTwiceDifferentiable}
% note: we don't formulate this result in Lean separately from \ref{convergence-for-twice-contdiff}
    There exists some $N_0 \in \N$ such that for all $N>N_0$ and $x\in [0,2\pi]$ we have
    \begin{equation}
        |S_N f_0 (x)- f_0(x)|\le \frac \epsilon 4\, .
    \end{equation}
\end{lemma}

We prove in \Cref{10difference}:
\begin{lemma}[control approximation effect]
\label{control-approximation-effect}
\uses{partial-Fourier-sums-of-small}
\leanok
\lean{control_approximation_effect}
    There is a set $E \subset \R$ with Lebesgue measure
    $|E|\le \epsilon$ such that for all
    \begin{equation}
        x\in [0,2\pi)\setminus E
    \end{equation}
   we have
   \begin{equation}
    \label{eq-max-partial-sum-diff}
    \sup_{N\ge 0} |S_Nf(x)-S_Nf_0(x)| \le \frac \epsilon 4\,.
    \end{equation}
\end{lemma}

We are now ready to prove \Cref{classical-carleson}. We first prove a version with explicit exceptional sets.

\begin{theorem}[classical Carleson with exceptional sets]
    \label{exceptional-set-carleson}
    \leanok
    \lean{exceptional_set_carleson}
    \uses{smooth-approximation, convergence-for-smooth,
    control-approximation-effect}
    Let $f$ be a $2\pi$-periodic complex-valued continuous function on $\mathbb{R}$. For all $\epsilon>0$, there exists a Borel set $E\subset [0,2\pi]$ with Lebesgue measure $|E|\le \epsilon$ and a positive integer $N_0$ such that for all $x\in [0,2\pi]\setminus E$ and all integers $N>N_0$, we have
    \begin{equation}\label{aeconv}
    |f(x)-S_N f(x)|\le \epsilon.
    \end{equation}
\end{theorem}
\begin{proof}
\leanok
Let $N_0$ be as in \Cref{convergence-for-smooth}.
For every
\begin{equation}
x\in [0, 2\pi) \setminus E\, ,
\end{equation}
and every $N>N_0$ we have by the triangle inequality
\begin{equation*}
    |f(x)-S_Nf(x)|
    \end{equation*}
    \begin{equation}\label{epsilonthird}
    \le |f(x)-f_0(x)|+ |f_0(x)-S_Nf_0(x)|+|S_Nf_0(x)-S_N f(x)|\, .
\end{equation}
Using \Cref{smooth-approximation,convergence-for-smooth,control-approximation-effect}, we estimate \eqref{epsilonthird} by
\begin{equation}
    \le \epsilon' +\frac \epsilon 4 +\frac \epsilon 4\le \epsilon\, .
\end{equation}
This shows \eqref{aeconv} for the given $E$ and $N_0$.
\end{proof}

Now we turn to the proof of the statement of Carleson theorem given in the introduction.


Let $\kappa:\R\to \C$ be the function defined by
$\kappa(0)=0$ and for $0<|x|<1$
\begin{equation}\label{eq-hilker}
\kappa(x)=\frac { 1-|x|}{1-e^{ix}}\,
\end{equation}
and for $|x|\ge 1$,
\begin{equation}\label{eq-hilker1}
\kappa(x)=0\, .
\end{equation}
Note that this function is continuous at every point $x$ with $|x|>0$.

The proof of \Cref{control-approximation-effect} will
use the following \Cref{real-Carleson}, which itself is proven
 in \Cref{10carleson} as an application of
 \Cref{metric-space-Carleson}.
\begin{lemma}[real Carleson]\label{real-Carleson}
\uses{two-sided-metric-space-Carleson,real-line-doubling,real-line-metric,real-line-measure,oscillation-control,frequency-monotone,frequency-ball-doubling,frequency-ball-growth,integer-ball-cover,real-van-der-Corput,Hilbert-strong-2-2,Hilbert-kernel-regularity,Hilbert-kernel-bound}
\leanok
\lean{rcarleson}
    Let $F,G$ be Borel subsets of $\R$ with finite measure. Let $f$ be a bounded measurable function on $\R$ with $|f|\le \mathbf{1}_F$. Then
\begin{equation}
    \left|\int _G Tf(x) \, dx\right| \le C_{4,2} |F|^{\frac 12} |G|^{\frac 12} \, ,
\end{equation}
where
\begin{equation}
    \label{define-T-carleson}
    T f(x)=\sup_{n\in \mathbb{Z}}
    \sup_{r>0}\left|\int_{r<|x-y|<1} f(y)\kappa(x-y) e^{iny}\, dy\right|\, .
\end{equation}
\end{lemma}


One of the main assumptions of \Cref{two-sided-metric-space-Carleson}, concerning the operator $T_r$ defined in \eqref{def-T-r}, is verified by the following lemma, which is proved in \Cref{10hilbert}.
\begin{lemma}[Hilbert strong 2 2]
    \label{Hilbert-strong-2-2}
    \uses{modulated-averaged-projection,integrable-bump-convolution,Dirichlet-approximation}
    \leanok
    \lean{Hilbert_strong_2_2}
    Let $0<r$. Let $f$ be a bounded, measurable function on $\mathbb{R}$. Then
    \begin{equation}
        \label{eq-Hr-L2-bound}
        \|H_rf\|_{2}\leq 2^{9} \|f\|_2,
    \end{equation}
    where
    \begin{equation}
        \label{def-H-r}
        H_r f(x) := T_r f(x) = \int_{r\le |x-y|} \kappa(x-y) f(y) \, dy
    \end{equation}
\end{lemma}

The next lemma will be used to verify that the collection $\Mf$ of modulation functions in our application of \Cref{metric-space-Carleson} satisfies the condition \eqref{eq-vdc-cond}.
It is proved in \Cref{10vandercorput}.

\begin{lemma}[van der Corput]
\label{van-der-Corput}
\leanok
\lean{van_der_Corput}
    Let $\alpha\le\beta$ be real numbers. Let $g:\R\to \C$ be a measurable function and assume
    \begin{equation}
        \|g\|_{Lip(\alpha,\beta)}:=\sup_{\alpha\le x\le \beta}|g(x)|+\frac{|\beta-\alpha|}{2}
        \sup_{\alpha\le x<y\le \beta} \frac {|g(y)-g(x)|}{|y-x|}<\infty\, .
    \end{equation}
    Then for any $\alpha \le \beta$ and $n\in\Z$ we have
    \begin{equation}
        \int _{\alpha}^{\beta} g(x) e^{inx}\, dx\le 2\pi |\beta-\alpha|\|g\|_{Lip(\alpha,\beta)}(1+|n||\beta-\alpha|)^{-1}\, .
    \end{equation}

\end{lemma}


We close this section with five lemmas that are used
across the following subsections.

\begin{lemma}[Dirichlet kernel]
\label{Dirichlet-kernel}
\leanok
\lean{dirichletKernel_eq, partialFourierSum_eq_conv_dirichletKernel}
We have for every $2\pi$-periodic bounded measurable $f$ and every $N\ge 0$
\begin{equation}
    S_Nf(x)=\frac 1{2\pi}\int_{0}^{2\pi}f(y) K_N(x-y)\, dy
\end{equation}
where $K_N$ is the $2\pi$-periodic continuous function of
$\R$ given by
\begin{equation}\label{eqksumexp}
\sum_{n=-N}^N e^{in x'}\, .
\end{equation}
We have for $e^{ix'}\neq 1$ that
\begin{equation}\label{eqksumhil}
    K_N(x')=\frac{e^{iNx'}}{1-e^{-ix'}}
      +\frac {e^{-iNx'}}{1-e^{ix'}} \, .
\end{equation}


\end{lemma}


\begin{proof}
\leanok
We have by definitions and interchanging sum and integral
   \begin{equation*}
        S_Nf(x)=\sum_{n=-N}^N \widehat{f}_n e^{inx}
    \end{equation*}
       \begin{equation*}
    =\sum_{n=-N}^N \frac 1{2\pi}\int_{0}^{2\pi}
    f(x) e^{in(x-y)}\, dy
    \end{equation*}
 \begin{equation}\label{eq-expsum}
     =\frac 1{2\pi}\int_{0}^{2\pi}
    f(y) \sum_{n=-N}^N e^{in(x-y)}\, dy\, .
 \end{equation}
This proves the first statement of the lemma.
By a telescoping sum, we have for every $x'\in \R$
\begin{equation}
    \left( e^{\frac 12 ix'}-e^{-\frac 12 ix'}\right) \sum_{n=-N}^N e^{inx'}= e^{(N+\frac 12) ix'}-e^{-(N+\frac 12) ix'}\, .
\end{equation}
If $e^{ix'}\neq 1$, the first factor on the left-hand side is not $0$ and we may divide by this factor to obtain
\begin{equation}
      \sum_{n=-N}^N e^{inx'}= \frac{e^{i(N+\frac 1 2)x'}}{e^{\frac 12 ix'}-e^{-\frac 12ix'}}
      -\frac{e^{-i(N+\frac 1 2)x'}}{e^{\frac 12 ix'}-e^{-\frac 12ix'}}
       =\frac{e^{iNx'}}{1-e^{-ix'}}
      +\frac {e^{-iNx'}}{1-e^{ix'}}\, .
\end{equation}
This proves the second part of the lemma.
\end{proof}

\begin{lemma}[lower secant bound]
    \label{lower-secant-bound}
    \leanok
    \lean{lower_secant_bound'}
    Let $\eta>0$ and $-2\pi +\eta \le x\le 2\pi-\eta$ with $|x|\ge \eta$. Then
    \begin{equation}
        |1-e^{ix}|\ge \frac{2}{\pi} \eta
    \end{equation}
\end{lemma}
\begin{proof}
    \leanok
    We have
    $$
        |1 - e^{ix}| = \sqrt{(1 - \cos(x))^2 + \sin^2(x)} \ge |\sin(x)|\,.
    $$
    If $0 \le x \le \frac{\pi}{2}$, then we have from concavity of $\sin$ on $[0, \pi]$ and $\sin(0) = 0$ and $\sin(\frac{\pi}{2}) = 1$
    $$
        |\sin(x)| \ge \frac{2}{\pi} x \ge \frac{2}{\pi} \eta\,.
    $$
    When $x\in \frac{m\pi}{2} + [0, \frac{\pi}{2}]$ for $m \in \{-4, -3, -2, -1, 1, 2, 3\}$ one can argue similarly.
\end{proof}

The following lemma will be proved in \Cref{10projection}.

\begin{lemma}[spectral projection bound]
\label{spectral-projection-bound}
\leanok
\lean{spectral_projection_bound}
    Let $f$ be a bounded $2\pi$-periodic measurable function. Then, for all $N\ge 0$
   \begin{equation}\label{snbound}
   \|S_Nf\|_{L^2[0, 2\pi]} \le \|f\|_{L^2[0, 2\pi]}.
   \end{equation}
\end{lemma}

\begin{lemma}[Hilbert kernel bound]
\label{Hilbert-kernel-bound}
\leanok
\lean{Hilbert_kernel_bound}
\uses{lower-secant-bound}
    For $x,y\in \R$ with $x\neq y$ we have
    \begin{equation}\label{eqcarl30}
        |\kappa(x-y)|\le 2^2(2|x-y|)^{-1}\, .
    \end{equation}
\end{lemma}
\begin{proof}
\leanok
    Fix $x\neq y$. If $\kappa(x-y)$ is zero, then \eqref{eqcarl30} is evident. Assume $\kappa(x-y)$ is not zero, then $0<|x-y|<1$.
    We have
\begin{equation}\label{eqcarl31}
|\kappa(x-y)|=\left|\frac {1-|x-y|}{1-e^{i(x-y)}}\right|\, .
\end{equation}
We estimate
with \Cref{lower-secant-bound}
\begin{equation}\label{eqcarl311}
|\kappa(x-y)|\le \frac {1}{|1-e^{i(x-y)}|}\le \frac 2{|x-y|}\, .
\end{equation}
This proves \eqref{eqcarl30} in the given case and completes the proof of the lemma.
\end{proof}

\begin{lemma}[Hilbert kernel regularity]
\label{Hilbert-kernel-regularity}
\leanok
\lean{Hilbert_kernel_regularity}
\uses{lower-secant-bound}
    For $x,y,y'\in \R$ with $x\neq y,y'$ and
    \begin{equation}
        \label{eq-close-hoelder}
        2|y-y'|\le |x-y|\, ,
    \end{equation}
    we have
    \begin{equation}\label{eqcarl301}
        |\kappa(x-y) - \kappa(x-y')|\le 2^{8}\frac{1}{|x-y|} \frac{|y-y'|}{|x-y|}\, .
    \end{equation}
\end{lemma}

\begin{proof}
\leanok
    Upon replacing $y$ by $y-x$ and $y'$ by $y'-x$ on the left-hand side of \eqref{eq-close-hoelder}, we can assume that $x = 0$. Then the assumption \eqref{eq-close-hoelder} implies that $y$ and $y'$ have the same sign. Since $\kappa(y) = \bar \kappa(-y)$ we can assume that they are both positive. Then it follows from \eqref{eq-close-hoelder} that
    $$
        \frac{y}{2} \le y' \,.
    $$
    We distinguish four cases. If $y, y' \le 1$, then we have
    $$
        |\kappa(-y) - \kappa(-y')| = \left| \frac{1 - y}{1- e^{-iy}} - \frac{1 - y'}{1- e^{-iy'}}\right|
    $$
    and by the fundamental theorem of calculus
    $$
        = \left| \int_{y'}^{y} \frac{-1 + e^{-it} + i(1-t)e^{it}}{(1 - e^{-it})^2} \,dt \right|\,.
    $$
    Using $y' \ge \frac{y}{2}$ and \Cref{lower-secant-bound}, we bound this by
    $$
        \le |y - y'| \sup_{\frac{y}{2} \le t \le 1} \frac{3}{|1 - e^{-it}|^2} \le 3 |y-y'| (2 \frac{2}{y})^2 \le 2^{6} \frac{|y-y'|}{|y|^2}\,.
    $$
    If $y \le 1$ and $y' > 1$, then $\kappa(-y') = 0$ and we have from the first case
    $$
        |\kappa(-y) - \kappa(-y')| = |\kappa(-y) - \kappa(-1)| \le 2^{6} \frac{|y-1|}{|y|^2} \le 2^{6} \frac{|y-y'|}{|y|^2}\,.
    $$
    Similarly, if $y > 1$ and $y' \le 1$, then $\kappa(-y) = 0$ and we have from the first case
    $$
        |\kappa(-y) - \kappa(-y')| = |\kappa(-y') - \kappa(-1)| \le 2^{6} \frac{|y'-1|}{|y'|^2} \le 2^{6} \frac{|y-y'|}{|y'|^2}\,.
    $$
    Using again $y' \ge \frac{y}{2}$, we bound this by
    $$
        \le 2^{6} \frac{|y-y'|}{|y / 2|^2} = 2^{8} \frac{|y-y'|}{|y|^2}
    $$
    Finally, if $y, y' > 1$ then
    $$
        |\kappa(-y) - \kappa(-y')| = 0 \le 2^{8} \frac{|y-y'|}{|y|^2}\,.
    $$
\end{proof}






\subsection{Smooth functions.}
\label{10smooth}
\begin{lemma}
\label{fourier-coeff-derivative}
\leanok
\lean{fourierCoeffOn_of_hasDerivAt}
Let $f:\R \to \C$ be $2\pi$-periodic and continuously differentiable, and let $n \in \Z \setminus \{0\}$. Then
\begin{equation}
    \widehat{f}_n = \frac{1}{i n} \widehat{f'}_n.
\end{equation}
\end{lemma}
\begin{proof}
\leanok
This is part of the Lean library.
\end{proof}

\begin{lemma}
\label{convergence-of-coeffs-summable}
\leanok
\lean{hasSum_fourier_series_of_summable}
Let $f:\R \to \C$ such that
\begin{equation}
    \sum_{n\in \Z} |\widehat{f}_n| < \infty.
\end{equation}
Then
\begin{equation}
    \sup_{x\in [0,2\pi]} |f(x) - S_Nf(x)| \rightarrow 0
\end{equation}
as $N \rightarrow \infty$.
\end{lemma}

\begin{proof}
\leanok
    This is part of the Lean library.
\end{proof}

\begin{lemma}
\label{convergence-for-twice-contdiff}
\uses{fourier-coeff-derivative,convergence-of-coeffs-summable}
\leanok
\lean{fourierConv_ofTwiceDifferentiable}
    Let $f:\R \to \C$ be $2\pi$-periodic and twice continuously differentiable. Then
    \begin{equation}
        \sup_{x\in [0,2\pi]} |f(x) - S_Nf(x)| \rightarrow 0
    \end{equation}
    as $N \rightarrow \infty$.
\end{lemma}
\begin{proof}
\leanok
By \Cref{convergence-of-coeffs-summable}, it suffices to show that the Fourier coefficients $\widehat{f}_n$ are summable.
Applying \Cref{fourier-coeff-derivative} twice and using the fact that $f''$ is continuous and thus bounded on $[0,2\pi]$ , we compute
\begin{equation*}
    \sum_{n\in \Z} |\widehat{f}_n| = |\widehat{f}_0| + \sum_{n\in \Z \setminus \{0\}} \frac {1}{n^2} |\widehat{f''}_n|
    \le |\widehat{f}_0| + \left(\sup_{x\in [0,2\pi]} |f(x)| \right) \cdot \sum_{n\in \Z \setminus \{0\}} \frac {1}{n^2}
    < \infty.
\end{equation*}
\end{proof}

\begin{proof}[Proof of \Cref{convergence-for-smooth}]
\leanok
\proves{convergence-for-smooth}
    \Cref{convergence-for-smooth} now follows directly from the previous \Cref{convergence-for-twice-contdiff}.
\end{proof}


\subsection{The truncated Hilbert transform}
\label{10hilbert}





Let $M_n$ be the modulation operator
acting on measurable $2\pi$-periodic functions
defined by
\begin{equation}
    M_ng(x)=g(x) e^{inx}\, .
\end{equation}
Define the approximate Hilbert transform by
\begin{equation}
    L_N g=\frac 1N\sum_{n=0}^{N-1}
       M_{n+N} S_{N+n}M_{-N-n}g\, .
\end{equation}


\begin{lemma}[modulated averaged projection]
\label{modulated-averaged-projection}
\leanok
\lean{modulated_averaged_projection}
\uses{spectral-projection-bound}
We have for every bounded measurable $2\pi$-periodic function $g$
\begin{equation}\label{lnbound}
    \|L_Ng\|_{L^2[0, 2\pi]}\le \|g\|_{L^2[0, 2\pi]}\,.
\end{equation}
\end{lemma}
\begin{proof}
    \leanok
    We have
    \begin{equation}\label{mnbound}
        \|M_ng\|_{L^2[0, 2\pi]}^2=\int_{0}^{2\pi} |e^{inx}g(x)|^2\, dx
        =\int_{0}^{2\pi} |g(x)|^2\, dx=\|g\|_{L^2[0, 2\pi]}^2\, .
    \end{equation}
     We have by the triangle inequality, the square root of the identity in \eqref{mnbound}, and \Cref{spectral-projection-bound}
    \begin{equation*}
        \|L_ng\|_{L^2[0, 2\pi]}=\|\frac 1N\sum_{n=0}^{N-1}
       M_{n+N} S_{N+n}M_{-N-n}g\|_{L^2[0, 2\pi]}
    \end{equation*}
    \begin{equation*}
        \le \frac 1N\sum_{n=0}^{N-1} \|
       M_{n+N} S_{N+n}M_{-N-n}g\|_{L^2[0, 2\pi]}
         = \frac 1N\sum_{n=0}^{N-1} \|
    S_{N+n}M_{-N-n}g\|_{L^2[0, 2\pi]}
    \end{equation*}
     \begin{equation}
     \le \frac 1N\sum_{n=0}^{N-1} \|
 M_{-N-n}g\|_{L^2[0, 2\pi]} = \frac 1N\sum_{n=0}^{N-1} \|
g\|_{L^2[0, 2\pi]} =\|g\|_{L^2[0, 2\pi]}\, .
    \end{equation}
This proves \eqref{lnbound} and completes the proof of the lemma.
\end{proof}

\begin{lemma}[periodic domain shift]\label{periodic-domain-shift}
\leanok
\lean{Function.Periodic.intervalIntegral_add_eq, intervalIntegral.integral_comp_sub_right}
Let $f$ be a bounded $2\pi$-periodic function. We have for any
$0 \le x\le 2\pi$ that
\begin{equation}
 \int_0^{2\pi} f(y)\, dy= \int_{-x}^{2\pi -x} f(y)\, dy
 =\int_{-\pi}^{\pi} f(y-x)\, dy\,.
\end{equation}
\end{lemma}
\begin{proof}
\leanok
    We have by periodicity and change of variables
    \begin{equation}\label{eqhil9}
 \int_{-x}^{0} f(y)\, dy=\int_{-x}^{0} f(y+2\pi)\, dy= \int_{2\pi -x}^{2\pi} f(y)\, dy\, .
\end{equation}
We then have by breaking up the domain of integration
and using \eqref{eqhil9}
\begin{equation*}
 \int_0^{2\pi} f(y)\, dy= \int_0^{2\pi -x} f(y)\, dy+
 \int_{2\pi -x}^{2\pi} f(y)\, dy
 \end{equation*}
\begin{equation}
= \int_0^{2\pi -x} f(y)\, dy+
 \int_{ -x}^{0} f(y)\, dy
 = \int_{-x}^{2\pi-x} f(y)\, dy\, .
 \end{equation}
This proves the first identity of the lemma. The second identity follows by substitution of $y$ by $y-x$.
\end{proof}



\begin{lemma}[Young convolution]\label{Young-convolution}
\leanok
\lean{young_convolution}
\uses{periodic-domain-shift}
    Let $f$ and $g$ be two bounded non-negative measurable $2\pi$-periodic functions on $\R$. Then
    \begin{equation}\label{eqyoung}
        \left(\int_0^{2\pi} \left(\int_0^{2\pi}
        f(y)g(x-y)\, dy\right)^2\, dx\right)^{\frac 12}\le \|f\|_{L^2[0, 2\pi]} \|g\|_{L^1[0, 2\pi]}\, .
    \end{equation}
    \end{lemma}
\begin{proof}
\leanok
    Using Fubini and \Cref{periodic-domain-shift}, we observe
\begin{equation*}
  \int_0^{2\pi}\int_0^{2\pi}f(y)^2g(x-y)\, dy
    \, dx=\int_0^{2\pi}f(y)^2\int_0^{2\pi}g(x-y)\, dx
    \, dy
\end{equation*}
\begin{equation}\label{eqhil4}
=\int_0^{2\pi}f(y)^2\int_0^{2\pi}g(x) \, dx
     dy
=\|f\|_{L^2[0, 2\pi]}^2\|g\|_{L^1[0, 2\pi]}\, .
\end{equation}
   Let $h$ be the nonnegative square root of $g$, then
   $h$ is bounded and $2\pi$-periodic with $h^2=g$.
   We estimate the square of the left-hand side of
   \eqref{eqyoung} with Cauchy-Schwarz and then with
   \eqref{eqhil4} by
       \begin{equation*}
         \int_0^{2\pi} (\int_0^{2\pi}f(y)h(x-y)h(x-y)\, dy)^2\, dx
   \end{equation*}
\begin{equation*}
    \le \int_0^{2\pi}\left(\int_0^{2\pi}f(y)^2g(x-y)\, dy\right)
    \left(\int_0^{2\pi}g(x-y)\, dy\right)\, dx
\end{equation*}
\begin{equation*}
    = \|f\|_{L^2[0, 2\pi]}^2\|g\|_{L^1[0, 2\pi]}^2\, .
\end{equation*}
Taking square roots, this proves the lemma.
\end{proof}

For $0<r<1$, Define the kernel $k_r$ to be the $2\pi$-periodic function
\begin{equation}
    k_r(x):=\min \left(r^{-1}, 1+\frac r{|1-e^{ix}|^2}\right)\, ,
\end{equation}
where the minimum is understood to be $r^{-1}$ in case $1=e^{ix}$.
\begin{lemma}[integrable bump convolution]
\label{integrable-bump-convolution}
\leanok
\lean{integrable_bump_convolution}
\uses{Young-convolution}
Let $g,f$ be bounded measurable $2\pi$-periodic functions. Let $0<r<\pi$.
Assume we have for all $x$
\begin{equation}\label{ebump1}
    |g(x)|\le k_r(x)\, .
\end{equation}
Let
\begin{equation}
   h(x)= \int_0^{2\pi} f(y)g(x-y)\, dy \, .
\end{equation}
Then
\begin{equation}
   \|h\|_{L^2[0, 2\pi]}\le 17\|f\|_{L^2[-\pi, \pi]} \, .
\end{equation}

\end{lemma}

\begin{proof}
\leanok
From monotonicity of the integral and \eqref{ebump1},
\begin{equation}
    \|g\|_{L^1[0, 2\pi]} \le \int_0^{2\pi}k_r(x)\, dx\,.
\end{equation}
Using the symmetry
$k_r(x)=k_r(-x)$, the assumption, and \Cref{lower-secant-bound}, the last display
is equal to
\begin{equation*}
    = 2 \int_0^\pi \min\left(\frac 1r, 1+\frac r{|1-e^{ix}|^2}\right)\, dx
\end{equation*}
\begin{equation*}
    \le 2\int_0^{r} \frac 1r \, dx+2\int_r^{\pi}1+\frac {64r}{x^2}\, dx
\end{equation*}
\begin{equation}
    \le 2+2\pi + 2\left(\frac {4r}r-\frac {4r}{\pi}\right)
    \le 17\, .
\end{equation}
    Together with \Cref{Young-convolution}, this proves the lemma.
\end{proof}

\begin{lemma}[Dirichlet approximation]\label{Dirichlet-approximation}
\uses{Dirichlet-kernel,lower-secant-bound}
\leanok
\lean{continuous_dirichletApprox, periodic_dirichletApprox, approxHilbertTransform_eq_dirichletApprox, dist_dirichletApprox_le}
Let $0<r<1$. Let $N$ be the smallest
integer larger than $\frac 1r$.
There is a $2\pi$-periodic continuous function
 ${L'}$ on $\R$ that satisfies for all $0\le x\le 2\pi$
and all $2\pi$-periodic bounded measurable functions $f$ on $\R$
\begin{equation}\label{lthroughlprime}
    L_Nf(x)=\frac 1{2\pi}\int_{0}^{2\pi}f(y) {L'}(x-y)\, dy\,.
\end{equation}
Moreover, for all $-\pi \le x \le \pi$,
\begin{equation}\label{eqdifflhil}
    \left|L'(x)-\mathbf{1}_{\{y:\, r<|y|<1\}} \kappa(x)\right|\le 12 k_r(x)\, .
\end{equation}
\end{lemma}


\begin{proof}
\leanok
We have by definition and \Cref{Dirichlet-kernel}
\begin{equation}
    L_Ng(x)=
    \frac 1N\sum_{n=0}^{N-1}
       \int_0^{2\pi} e^{-i(N+n)x} K_{N+n}(x-y) e^{i(N+n)y}g(y)
\, dy \, .\end{equation}
This is of the form \eqref{lthroughlprime} with
the continuous function
\begin{equation}
    {L'}(x)= \frac 1N\sum_{n=0}^{N-1}
      K_{N+n}(x) e^{-i(N+n)x}\, .
\end{equation}
With \eqref{eqksumexp} of \Cref{Dirichlet-kernel}
we have $|K_N(x)|\le 2N+1$ for every $x$ and thus
\begin{equation}\label{eqhil13}
    |{L'}(x)|\le \frac 1N\sum_{n=0}^{N-1}
      (2N+2n+1) \le 4N\le 2^3 r^{-1}\, .
\end{equation}
Therefore, for $|x|\in [0, r)\cup (1, \pi]$, we have
\begin{equation}
    \label{eqdiffzero}
    \left|L'(x)-\mathbf{1}_{\{y:\, r<|y|<1\}}(x)\kappa(x)\right|=|L'(x)|\leq 2^{3} r^{-1}.
\end{equation}
This proves \eqref{eqdifflhil} for $|x|\in [0, r)$ since $k_r(x)=r^{-1}$ in this case.

For $e^{ix}\neq 1$
one may use the expression
\eqref{eqksumhil} for $K_N$
in \Cref{Dirichlet-kernel} to obtain
\begin{equation*}
    {L'}(x)= \frac 1N\sum_{n=0}^{N-1}
     \left(\frac{e^{i(N+n)x}}{1-e^{-ix}}
      +\frac {e^{-i(N+n)x}}{1-e^{ix}}\right) e^{i(N+n)x}
\end{equation*}
\begin{equation*}
    = \frac 1N\sum_{n=0}^{N-1}
    \left(\frac{e^{i2(N+n)x}}{1-e^{-ix}}
      +\frac {1}{1-e^{ix}}\right)
\end{equation*}
\begin{equation}\label{eqhil3}
    = \frac{1}{1-e^{ix}} +
     \frac 1N \frac {e^{i2Nx}}{1-e^{-ix}}
     \sum_{n=0}^{N-1}
    {e^{i2nx}}
\end{equation}
and thus
\begin{equation}
\label{eq-L'L''}
  {L'}(x) -\mathbf{1}_{\{y:\, r<|y|<1\}}\kappa(x)=L''(x)+ \frac{1-\mathbf{1}_{{\{y:\, r<|y|<1\}}}(x)(1-|x|)}{1-e^{ix}},
\end{equation}
where
$$L''(x):=\frac 1N \frac {e^{i2Nx}}{1-e^{-ix}}
     \sum_{n=0}^{N-1}
    {e^{i2nx}}.$$
For $x\in [-\pi, r]\cup [r, \pi]$, we have using \Cref{lower-secant-bound} that
\begin{equation*}
   \left|\frac{1-\mathbf{1}_{{\{y:\, r<|y|<1\}}}(x)(1-|x|)}{1-e^{ix}} \right|=\left|\frac{\min(|x|, 1)}{1-e^{ix}} \right|\leq \frac{2\min(|x|, 1)}{|x|}
\end{equation*}
\begin{equation}
 \label{eq-diffzero2}
    \leq 2\cdot 1\leq 2 k_r(x).
\end{equation}
Next, we need to estimate $L''(x)$. If the real part of
$e^{ix}$ is negative, we have
\begin{equation}
  1\le |1-e^{-ix}|\le 2\, .
\end{equation}
and hence
\begin{equation}\label{eqhil12}
    |L''(x)|\le
     \frac 1N
     \sum_{n=0}^{N-1}
    1=1\le 1+\frac r{|1-e^{ix}|^2}\, .
\end{equation}
If the real part of $e^{ix}$ is positive and in particular while still $e^{ix}\neq \pm 1$, then we have by telescoping
\begin{equation}
 (1-e^{2ix})
     \sum_{n=0}^{N-1}
    {e^{i2nx}}=1-e^{i2Nx}\, .
\end{equation}
As $e^{2ix}\neq 1$, we may divide by $1-e^{2ix}$ and insert this into
\eqref{eqhil3} to obtain
\begin{equation}
 L''(x)=
           \frac 1N \frac {e^{i2Nx}}{1-e^{-ix}}
     \frac{1-e^{i2Nx}}{1-e^{2ix}}\, .
\end{equation}
Hence, using the nonnegativity of the real part of $e^{ix}$
\begin{equation*}
    |L''(x)|
 \le \frac 2 N \frac {1}{|1-e^{ix}|}
     \frac{1}{|1-e^{2ix}|}
 \end{equation*}
\begin{equation}\label{eqhil11}
    = \frac 2 N \frac {1}{|1-e^{ix}|^2}
     \frac{1}{|1+e^{ix}|}\le
 \frac {2r}{|1-e^{ix}|^2}\le 2 \left(1+\frac {r}{|1-e^{ix}|^2}\right)
\end{equation}
Moreover, for $|x| \in [r, \pi]$, one checks easily that
\begin{equation*}
  1+\frac {r}{|1-e^{ix}|^2} \leq 5 k_r(x).
\end{equation*}
Therefore,~\eqref{eqhil12} and~\eqref{eqhil11} show that, in this range of $x$,
\begin{equation*}
  |L''(x)| \le 10 k_r(x).
\end{equation*}
Together with~\eqref{eq-diffzero2}, this prove \eqref{eqdifflhil} for $|x| \in [r, \pi]$.
As the range $|x| \in [0,r)$ is covered by~\eqref{eqdiffzero}, this completes the proof of the lemma.
\end{proof}


We now prove \Cref{Hilbert-strong-2-2}.

\begin{proof}[Proof of \Cref{Hilbert-strong-2-2}]
    \proves{Hilbert-strong-2-2}
    \leanok
    We first show that if $f$ is supported in $[1, 4]$, then
    \begin{equation}
        \label{eq-Hr-short-support}
        \|H_r f\|_{L^2[2, 3]} \le 2^{8} \|f\|_{L^2(\R)}\,.
    \end{equation}
    Let $\tilde{f}$ be the $2\pi$-periodic extension of $f$ to $\mathbb{R}$. Let $N$ be the smallest
    integer larger than $\frac 1r$. Then, \Cref{eqdifflhil} shows that the kernels of $H_r$ and $2\pi L_N$ differ by at
    most $12k_r$ on $[-\pi, \pi]$. Consider $x \in [2,3]$. When computing $H_r f(x)$ and $2\pi L_N f(x)$, the kernels are
    computed at points of the form $x-y$ with $f(y) \ne 0$, i.e., $y \in [1,4]$. As $x\in [2,3]$, the difference
    $x-y$ is bounded in absolute value by $2$, and therefore belongs to $[-\pi, \pi]$, where the above bound holds.

    It follows that, for $x \in [2,3]$,
    \begin{equation*}
        |H_r \tilde{f}(x)|\leq 2\pi |L_N \tilde{f}(x)|+12\left|\int_{0}^{2\pi}k_r(x-y)\tilde{f}(y)\, dy\right|.
    \end{equation*}
    Taking $L^2$ norm and using its sub-additivity, we get
    \begin{equation*}
       \|H_r \tilde{f}\|_{L^2[2, 3]}\leq 2\pi \|L_N \tilde{f}\|_{L^2[0, 2\pi]}\, + 12\left(\int_{0}^{2\pi} \left|\int_{0}^{2\pi}k_r(x-y)\tilde{f}(y)\, dy\right|^2\, dx\right)^{\frac{1}{2}}.
    \end{equation*}
    Since $\kappa$ is supported in $[-1,1]$, we have that $H_r\tilde{f}$ agrees with $H_r f$ on $[2,3]$.
    Using \Cref{modulated-averaged-projection} and \Cref{integrable-bump-convolution}, we conclude
    \begin{equation}
        \|H_r f\|_{L^2[2, 3]} \le 2\pi \|f\|_{L^2[0, 2\pi]} + 12 \cdot 17 \|f\|_{L^2[0, 2\pi]}\,,
    \end{equation}
    which gives \eqref{eq-Hr-short-support}.

    Suppose now that $f$ is supported in $[c, c+3]$ for some $c \in \R$. Then the function $g(x) = f(x-c+1)$ is supported in $[1, 4]$.
    By a change of variables in \eqref{def-H-r}, we have $H_r g(x ) = H_r f(x - c + 1)$. Thus, by \eqref{eq-Hr-short-support}
    \begin{equation}
        \label{eq-Hr-short-support-2}
        \|H_rf\|_{L^2[c+1, c+2]} = \|H_r g\|_{L^2[2,3]} \le 2^{8} \|g\|_2 = 2^8 \|f\|_2\,.
    \end{equation}

    Let now $f$ be arbitrary.
    Since $\kappa(x) = 0$ for $|x| > 1$, we have for all $x \in [c+1, c+2]$
    $$
        H_rf(x) = H_r(f \mathbf{1}_{[c, c+3]})(x)\,.
    $$
    Thus
    $$
        \int_{c+1}^{c+2} |H_r f(x)|^2 \, \mathrm{d}x = \int_{c+1}^{c+2}|H_r(f \mathbf{1}_{[c, c+3]})(x)|^2 \, \mathrm{d}x\,.
    $$
    Applying the bound \eqref{eq-Hr-short-support-2}, this is
    $$
        \le 2^{16} \int_{c}^{c+3} |f(x)|^2 \, \mathrm{d}x\,.
    $$
    Summing over all $c \in \mathbb{Z}$, we obtain
    $$
        \int_{\R} |H_rf(x)|^2 \, \mathrm{d}x \le 3 \cdot 2^{16} \int_{\R} |f(x)|^2 \, \mathrm{d}x\,.
    $$
    This completes the proof.
\end{proof}







\subsection{The proof of the van der Corput Lemma}
\label{10vandercorput}

\begin{proof}[Proof of \Cref{van-der-Corput}]
\proves{van-der-Corput}
\leanok
Let $g$ be a Lipschitz continuous function as in the lemma.
Assume first that $n=0$. Then
\begin{equation*}
    \int_\alpha^\beta g(x) \, \mathrm{d}x \le |\beta - \alpha| \sup_{\alpha\le x\le \beta}|g(x)|
    \le |\beta-\alpha|\|g\|_{Lip(\alpha,\beta)}(1+|n||\beta-\alpha|)^{-1}
\end{equation*}
Assume now $n\ne 0$. Without loss of generality, we may assume $n>0$.
We distinguish two cases. If $\beta-\alpha < \frac{\pi}{n}$, we have by the triangle inequality
\begin{equation*}
    \left|\int_\alpha^\beta g(x) e^{inx} \, \mathrm{d}x\right|
    \le |\beta -\alpha| \sup_{x \in [\alpha,\beta]} |g(x)|
    \le 2\pi |\beta-\alpha|\|g\|_{Lip(\alpha,\beta)}(1+|n||\beta-\alpha|)^{-1} \,.
\end{equation*}
We turn to the case $\frac{\pi}{n} \le \beta-\alpha$.
We have
$$
    e^{in(x + \pi/n)} = -e^{inx}\,.
$$
Using this, we write
$$
    \int_\alpha^\beta g(x) e^{inx} \, \mathrm{d}x
    = \frac{1}{2} \int_\alpha^\beta g(x) e^{inx} \, \mathrm{d}x - \frac{1}{2} \int_\alpha^\beta g(x) e^{in(x + \pi/n)}) \, \mathrm{d}x\,.
$$
We split the first integral at $\alpha + \frac{\pi}{n}$ and the second one at $\beta - \frac{\pi}{n}$, and make a change of variables in the second part of the first integral to obtain
$$
    = \frac{1}{2} \int_{\alpha}^{\alpha + \frac{\pi}{n}} g(x) e^{inx} \, \mathrm{d}x - \frac{1}{2} \int_{\beta - \frac{\pi}{n}}^{\beta} g(x) e^{in(x + \pi/n)} \, \mathrm{d}x
$$
$$
    + \frac{1}{2} \int_{\alpha + \frac{\pi}{n}}^{\beta} (g(x) - g(x - \frac{\pi}{n})) e^{inx} \, \mathrm{d}x\,.
$$
The sum of the first two terms is by the triangle inequality bounded by
$$
    \frac{\pi}{n} \sup_{x \in [\alpha,\beta]} |g(x)|\,.
$$
The third term is by the triangle inequality at most
$$
    \frac{1}{2} \int_{\alpha + \frac{\pi}{n}}^\beta |g(x) - g(x - \frac{\pi}{n})| \, \mathrm{d}x
$$
$$
    \le \frac{|\beta-\alpha|}{2} \frac{\pi}{n} \sup_{\alpha \le x < y \le \beta} \frac{|g(x) - g(y)|}{|x-y|}\,.
$$
Adding the two terms, we obtain
$$
    \left|\int_\alpha^\beta g(x) e^{-inx} \, \mathrm{d}x\right| \le \frac{\pi}{n} \|g\|_{\mathrm{Lip}(\alpha,\beta)}\,.
$$
This completes the proof of the lemma, using that with $\frac{\pi}{n} \le \beta-\alpha$,
$$
    \frac{\pi}{n} = \frac{2 \pi |\beta-\alpha|}{2n|\beta-\alpha|} \le 2 \pi |\beta-\alpha|(1 + n|\beta-\alpha|)^{-1}\,.
$$
\end{proof}


\subsection{Partial sums as orthogonal projections}
\label{10projection}

This subsection proves \Cref{spectral-projection-bound}

\begin{proof}[Proof of \Cref{spectral-projection-bound}]
\proves{spectral-projection-bound}
\leanok

The functions $e_n:x\mapsto e^{inx}$ form an orthonormal basis in $L^2[-\pi, \pi]$
(this is already in Mathlib).
Therefore we have
\begin{align*}
    \|S_Nf\|^2_{L^2[-\pi, \pi]}
    &= \|\sum_{n=-N}^N \langle\widehat{f}_n, e_n\rangle_{L^2[-\pi, \pi]}\|^2_{L^2[-\pi, \pi]} \\
    &= \sum_{n=-N}^N |\widehat{f}_n| \\
    &\le \sum_{n\in \mathbb{Z}} |\widehat{f}_n| \\
    &= \|f\|_{L^2[-\pi, \pi]}.
\end{align*}

This completes the proof of the lemma.
\end{proof}

\subsection{The error bound}
\label{10difference}

\begin{lemma}[Dirichlet kernel - Hilbert kernel relation]
    \label{Dirichlet-Hilbert}
    \leanok
    \lean{Dirichlet_Hilbert_diff}
    \uses{Dirichlet-kernel,lower-secant-bound}
    For all $N\in\Z$ and $x\in [-\pi,\pi] \setminus \{0\}$,
    \begin{equation*}
        \left|K_N(x) - (e^{-iNx}\kappa(x) + \overline{e^{-iNx}\kappa(x)})\right| \le \pi \,.
    \end{equation*}
\end{lemma}

\begin{proof}
    \leanok
    Let $N\in\Z$ and $x\in [-\pi,\pi] \setminus \{0\}$. With \Cref{Dirichlet-kernel}, we obtain
    \begin{equation*}
        K_N(x) - (e^{-iNx}\kappa(x) + \overline{e^{-iNx}\kappa(x)})
        = e^{-iNx} \frac{\min(|x|, 1) }{1 - e^{ix}} + e^{iNx} \frac{\min(|x|, 1) }{1 - e^{-ix}} \,.
    \end{equation*}
    Using \Cref{lower-secant-bound} with $\eta = \min(|x|, 1)$, we bound
    \begin{equation*}
        \left|K_N(x) - (e^{-iNx}\kappa(x) + \overline{e^{-iNx}\kappa(x)})\right|
        \le \frac{\min(|x|, 1) }{|1 - e^{ix}|} + \frac{\min(|x|, 1)}{|1 - e^{-ix}|}
        \le \frac{\pi}{2} + \frac{\pi}{2} = \pi \,.
    \end{equation*}
\end{proof}

\begin{lemma}[partial Fourier sum bound]
    \label{partial-Fourier-sum-bound}
    \uses{Dirichlet-kernel, Dirichlet-Hilbert}
    \leanok
    \lean{partialFourierSum_bound}
    Let $g:\R\to\C$ be a measurable $2\pi$-periodic function such that for some $\delta>0$ and every $x\in\R$,
    \begin{equation}
        |g(x)|\le \delta \,.
    \end{equation}
    Then for every $x\in [0,2\pi]$ and $N>0$,
    \begin{equation*}
        |S_N g(x)| \le \frac{1}{2\pi} (Tg(x) + T\bar{g}(x)) + \pi\delta.
    \end{equation*}
\end{lemma}

\begin{proof}
    \leanok
    Let $x\in [0,2\pi]$ and $N>0$. We have with \Cref{Dirichlet-kernel}
    \begin{equation*}
        |S_N g(x)| = \frac{1}{2\pi} \left| \int_0^{2\pi} g(y) K_N(x-y) \, dy\right|\,.
    \end{equation*}
    We use $2\pi$-periodicity of $g$ and $K_N$ to shift the domain of integration to obtain
    \begin{equation*}
        = \frac{1}{2\pi} \left|\int_{x-\pi}^{x+\pi} g(y) K_N(x-y) \, dy\right|\,.
    \end{equation*}
    Using the triangle inequality, we split this as
    \begin{equation}
        \label{eq-diff-integrable}
        \le \frac{1}{2\pi} \left|\int_{x-\pi}^{x+\pi} g(y) \left(K_N(x-y) - \max(|x-y|,0) K_N(x-y)\right) \, dy \right|
    \end{equation}
    \begin{equation}
        \label{eq-diff-singular}
        + \frac{1}{2\pi} \left|\int_{x-\pi}^{x+\pi} g(y) \max(|x-y|,0) K_N(x-y) \, dy\right|\,.
    \end{equation}
    Note that all integrals are well defined, since $K_N$ is by \eqref{eqksumexp} bounded by $2N+1$.
    Using that
    \begin{equation}
        \label{eq-Dirichlet-Hilbert}
        \max(|x-y|,0) K_N(x-y) = e^{-iN(x-y)}\kappa(x-y) + \overline{e^{-iN(x-y)}\kappa(x-y)} \,,
    \end{equation}
    \Cref{Dirichlet-Hilbert} and \eqref{g-small}, we bound \eqref{eq-diff-integrable} by
    \begin{equation*}
        \frac{1}{2\pi} \int_{x-\pi}^{x+\pi} |g(y)| \left|K_N(x-y) - e^{-iN(x-y)}\kappa(x-y) + \overline{e^{-iN(x-y)}\kappa(x-y)}\right|\, dy
        \le \pi\delta \,.
    \end{equation*}
    By dominated convergence and since $\kappa(x-y) = 0$ for $|x-y| > 1$, \eqref{eq-diff-singular} equals
    \begin{equation*}
        \frac{1}{2\pi} \lim_{r \to 0^+} \left| \int_{r < |x-y| < 1} g(y) \max(|x-y|,0) K_N(x-y) \, dy\right|\,.
    \end{equation*}
    We bound the limit by a supremum and rewrite using \eqref{eq-Dirichlet-Hilbert},
    \begin{equation*}
        \le \frac{1}{2\pi} \sup_{r > 0} \left| \int_{r < |x-y| < 1} g(y) \left(e^{-iN(x-y)}\kappa(x-y) + \overline{e^{-iN(x-y)}\kappa(x-y)}\right) \, dy\right|
    \end{equation*}
    Using the triangle inequality, we further bound this by
    \begin{alignat*}{3}
        \le&&&\frac{1}{2\pi} \sup_{r > 0} \left| \int_{r < |x-y| < 1} g(y) e^{-iNy} \kappa(x-y) \, dy\right| \\
        &+ &&\frac{1}{2\pi} \sup_{r > 0} \left| \int_{r < |x-y| < 1} \overline{g}(y) e^{-iNy} \kappa(x-y) \, dy\right|\,.
    \end{alignat*}
    By the definition \eqref{define-T-carleson} of $T$, this is
    \begin{equation*}
        \le \frac{1}{2\pi} (Tg(x) + T\bar{g}(x))\,.
    \end{equation*}
\end{proof}

%fixme: Put this somewhere else?
\begin{lemma}[real Carleson operator measurable]
    \label{real-Carleson-operator-measurable}
    \uses{Hilbert-kernel-bound}
    \leanok
    \lean{carlesonOperatorReal_measurable}
    Let $f$ be a bounded measurable function on $\R$. Then $Tf$ as defined in \eqref{define-T-carleson} is measurable.
\end{lemma}
\begin{proof}
    \leanok
    Since a countable supremum of measurable functions is measurable, it suffices to show that for every $n\in\Z$,
    \begin{equation*}
        x \mapsto \sup_{r>0}\left|\int_{r<|x-y|<1} f(y)\kappa(x-y) e^{iny}\, dy\right|
    \end{equation*}
    is measurable. So let $n\in\Z$.
    Note that for each $x\in\R$, the function
    \begin{equation*}
        r \mapsto \left|\int_{r<|x-y|<1} f(y)\kappa(x-y) e^{iny}\, dy\right|
    \end{equation*}
    is continuous on $(0,\infty)$ since the integrand is locally bounded on the domain $0<|x-y|<1$ by the assumptions on $f$ and \Cref{Hilbert-kernel-bound}.
    Thus, for each $x\in\R$,
    \begin{equation*}
        \sup_{r>0}\left|\int_{r<|x-y|<1} f(y)\kappa(x-y) e^{iny}\, dy\right|
        =\sup_{r\in\Q_{>0}}\left|\int_{r<|x-y|<1} f(y)\kappa(x-y) e^{iny}\, dy\right|
    \end{equation*}
    The right hand side is again a countable supremum so it remains to show that for every $r\in\Q_{>0}$,
    \begin{equation*}
        x \mapsto \left|\int_{r<|x-y|<1} f(y)\kappa(x-y) e^{iny}\, dy\right| = \left|\int \mathbf{1}_{\{r<|x-\cdot|<1\}}(y) f(y)\kappa(x-y) e^{iny}\, dy\right|
    \end{equation*}
    is measurable, which follows from the fact that the integrand is measurable in $(x,y)$.
\end{proof}

\begin{lemma}[partial Fourier sums of small]
    \label{partial-Fourier-sums-of-small}
    \uses{real-Carleson-operator-measurable,real-Carleson,partial-Fourier-sum-bound}
    \leanok
    \lean{control_approximation_effect} % todo: fix. This seems currently inlined in control_approximation_effect ?
    %fixme: rename \delta to \epsilon' ?
    Let $g:\R\to\C$ be a measurable $2\pi$-periodic function such that for some $\delta>0$ and every $x\in\R$,
    \begin{equation} \label{g-small}
        |g(x)|\le \delta \,.
    \end{equation}
    Then for every $\epsilon>0$, there exists a measurable set $E\subset [0,2\pi]$ with $|E|<\epsilon$ such that for every $x\in [0,2\pi]\setminus E$ and $N>0$,
    \begin{equation}  \label{S-Ng-small}
        |S_N g(x)|\le C_\epsilon \delta,
    \end{equation}
    where
    %fixme: replace by C(\epsilon) ?
    \begin{equation} \label{C-epsilon-def}
        C_\epsilon = \left(\frac{8}{\pi\epsilon}\right)^\frac{1}{2} C_{4,2} + \pi \,.
    \end{equation}
\end{lemma}

\begin{proof}
\leanok
Define
\begin{equation*}
    E := \{x \in [0, 2\pi] \ : \ \sup_{N > 0} |S_N g (x)| > C_\epsilon \delta \} \,.
\end{equation*}
Then \eqref{S-Ng-small} clearly holds, and it remains to show that $|E| \le \epsilon$.
With \Cref{partial-Fourier-sum-bound}, we obtain
\begin{equation*}
    E \subset \{x\in [0,2\pi] \ : \ C_\epsilon \delta < \frac{1}{2\pi} (Tg(x) + T\bar{g}(x)) + \pi\delta\} \subset E_1 \cup E_2,
\end{equation*}
where
\begin{align*}
    E_1 :=& \{x\in [0,2\pi] \ : \ \pi(C_\epsilon - \pi) \delta < Tg(x)\} \\
    E_2 :=& \{x\in [0,2\pi] \ : \ \pi(C_\epsilon - \pi) \delta < T\bar{g}(x)\}.
\end{align*}
By \Cref{real-Carleson-operator-measurable}, $E_1$ and $E_2$ are measurable. Thus,
\begin{equation*}
    \pi(C_\epsilon - \pi) \delta |E_1| \le \int_{E_1} Tg(x) \, dx = \delta \int_{E_1} T(\delta^{-1} g\mathbf{1}_{[-\pi,3\pi]})(x) \, dx \,.
\end{equation*}
Applying \Cref{real-Carleson} with $F = [-\pi, 3\pi]$ and $G = E'$, it follows that this is
\begin{equation*}
    \le \delta \cdot C_{4,2} |F|^{\frac{1}{2}} |E_1|^{\frac{1}{2}} \le (4\pi)^\frac{1}{2} C_{4,2} \delta \cdot |E'|^{\frac{1}{2}}\,.
\end{equation*}
Rearranging, we obtain
\begin{equation*}
    |E_1| \le \left(\frac{(4\pi)^\frac{1}{2} C_{4,2}}{\pi(C_\epsilon - \pi)}\right)^2 = \frac{\epsilon}{2}\,.
\end{equation*}
Analogously, we get the same estimate for $|E_2|$. This completes the proof using $|E| \le |E_1| + |E_2|$.
\end{proof}

%fixme: It seems unnecessary to have the extra lemma with this proof that is really only application.
\begin{proof}[Proof of \Cref{control-approximation-effect}]
\proves{control-approximation-effect}
\leanok
    \Cref{control-approximation-effect} follows directly from \Cref{partial-Fourier-sums-of-small} with $\delta:=\epsilon'$.
\end{proof}




\subsection{Carleson on the real line}
\label{10carleson}

We prove \Cref{real-Carleson}.

Consider the standard distance function
\begin{equation}
    \rho(x,y)=|x-y|
\end{equation}
on the real line $\R$.
\begin{lemma}[real line metric]
\label{real-line-metric}
\leanok
\lean{instProperSpaceReal, locallyCompact_of_proper, Real.instCompleteSpace}
The space $(\R,\rho)$ is a complete locally compact metric space.
\end{lemma}
\begin{proof}
    \leanok
    This is part of the Lean library.
\end{proof}
\begin{lemma}[real line ball]
\label{real-line-ball}
    \leanok
    \lean{Real.ball_eq_Ioo}
    For $x\in R$ and $R>0$, the ball $B(x,R)$ is the interval $(x-R,x+R)$
\end{lemma}
\begin{proof}
\leanok
Let $x'\in B(x,R)$. By definition of the ball,
$|x'-x|<R$. It follows that $x'-x<R$ and $x-x'<R$.
It follows $x'<x+R$ and $x'>x-R$. This implies
$x'\in (x-R,x+R)$.
Conversely, let $x'\in (x-R,x+R)$. Then
$x'<x+R$ and $x'>x-R$. It follows that
$x'-x<R$ and $x-x'<R$. It follows that $|x'-x|<R$,
hence $x'\in B(x,R)$.
This proves the lemma.
\end{proof}
We consider the Lebesgue measure $\mu$ on $\R$.
\begin{lemma}[real line measure]
\label{real-line-measure}
\leanok
\lean{instIsAddHaarMeasureVolume}
    The measure $\mu$ is a sigma-finite non-zero
    Radon-Borel measure on $\R$.
\end{lemma}
\begin{proof}
    \leanok
    This is part of the Lean library.
\end{proof}
\begin{lemma}[real line ball measure]
\label{real-line-ball-measure}
\uses{real-line-ball}
    \leanok
    \lean{Real.volume_ball}
    We have for every $x\in \R$ and $R>0$
    \begin{equation}
        \mu(B(x,R))=2R\, .
    \end{equation}
\end{lemma}
\begin{proof}
\leanok
We have with \Cref{real-line-ball}
\begin{equation}
    \mu(B(x,R))=\mu((x-R,x+R))=2R\, .
\end{equation}
\end{proof}

\begin{lemma}[real line doubling]
\label{real-line-doubling}
\leanok
\lean{MeasureTheory.InnerProductSpace.IsDoubling}
\uses{real-line-ball-measure}
    We have for every $x\in \R$ and $R>0$
    \begin{equation}
        \mu(B(x,2R))=2\mu(B(x,R))\, .
    \end{equation}
\end{lemma}
\begin{proof}
    \leanok
    We have with \Cref{real-line-ball-measure}
\begin{equation}
    \mu(B(x,2R)=4R=2\mu(B(x,R))\, .
\end{equation}
This proves the lemma.
\end{proof}


The preceding four lemmas show that $(\R, \rho, \mu, 4)$ is a doubling metric measure space. Indeed, we even show
that $(\R, \rho, \mu, 1)$ is a doubling metric measure space, but we may relax the estimate in \Cref{real-line-doubling} to conclude that $(\R, \rho, \mu, 4)$
is a doubling metric measure space.


For each $n\in \mathbb{Z}$ define
$\mfa_n:\R\to \R$ by
\begin{equation}
    \mfa_n(x)=nx\, .
\end{equation}


Let $\Mf$ be the collection $\{\mfa_n, n\in \mathbb{Z}\}$.
Note that for every $n\in \mathbb{Z}$ we have $\mfa_n(0)=0$.
Define
\begin{equation}
    \label{eqcarl4}
    d_{B(x,R)}(\mfa_n, \mfa_m) := 2R|n-m|\,.
\end{equation}

\begin{lemma}[frequency metric]
    \label{frequency-metric}
    \leanok
    \lean{instFunctionDistancesReal}
    For every $R > 0$ and $x \in X$, the function $d_{B(x,R)}$ is a metric on $\Mf$.
\end{lemma}

\begin{proof}
    \leanok
    This follows immediately from the fact that the standard metric on $\mathbb{Z}$ is a metric.
\end{proof}

\begin{lemma}[oscillation control]
\label{oscillation-control}
\leanok
\lean{oscillation_control}
    For every $R > 0$ and $x \in X$, and for all $n, m \in \mathbb{Z}$, we have
    \begin{equation}\label{eqcarl2}
       \sup_{y,y'\in B(x,R)}|ny-ny'-my+my'|\le 2|n-m|R\, .
    \end{equation}
\end{lemma}

\begin{proof}
\leanok
    The right hand side of \eqref{eqcarl2} equals
    $$
        \sup_{y,y'\in B(x,R)}|(n-m)(y-x)-(n-m)(y'-x)|\,.
    $$
    The lemma then follows from the triangle inequality.
\end{proof}

\begin{lemma}[frequency monotone]
\label{frequency-monotone}
\leanok
\lean{frequency_monotone}

    For any $x, x' \in X$ and $R, R' > 0$ with $B(x,R) \subset B(x, R')$, and for any $n, m \in \mathbb{Z}$
    $$
        d_{B(x,R)}(\mfa_n, \mfa_m) \le d_{B(x',R')}(\mfa_n, \mfa_m)\,.
    $$
\end{lemma}

\begin{proof}
\leanok
    This follows immediately from the definition \eqref{eqcarl4} and $R \le R'$.
\end{proof}

\begin{lemma}[frequency ball doubling]
\label{frequency-ball-doubling}
\leanok
\lean{frequency_ball_doubling}
\uses{frequency-metric}
  For any $x,x'\in \R$ and $R>0$ with
   $x\in B(x',2R)$ and any $n,m\in \mathbb{Z}$, we have
\begin{equation}\label{firstdb1}
    d_{B(x',2R)}(\mfa_n,\mfa_m)\le 2 d_{B(x,R)}(\mfa_n,\mfa_m) \, .
\end{equation}
\end{lemma}

\begin{proof}
\leanok
With \eqref{eqcarl4}, both sides of \eqref{firstdb1} are equal to $4R|n-m|$. This proves the lemma.
\end{proof}

\begin{lemma}[frequency ball growth]
\label{frequency-ball-growth}
\leanok
\lean{frequency_ball_growth}
    For any $x,x'\in \R$ and $R>0$ with
   $B(x,R)\subset B(x',2R)$ and any $n,m\in \mathbb{Z}$, we have
\begin{equation}\label{seconddb1}
    2d_{B(x,R)}(\mfa_n,\mfa_m)\le d_{B(x',2R)}(\mfa_n,\mfa_m) \, .
\end{equation}
\end{lemma}

\begin{proof}
\leanok
    With \eqref{eqcarl4}, both sides of \eqref{firstdb1} are equal to $4R|n-m|$. This proves the lemma.
\end{proof}

\begin{lemma}[integer ball cover]
    \label{integer-ball-cover}
    \leanok
    \lean{integer_ball_cover}
    For every $x\in \R$ and $R>0$ and every
    $n\in \mathbb{Z}$ and $R'>0$,
    there exist $m_1, m_2, m_3\in \mathbb{Z}$
    such that
    \begin{equation}\label{eqcarl5}
        B'\subset B_1\cup B_2\cup B_3\, ,
    \end{equation}
where
\begin{equation}
B'= \{ \mfa \in \Mf: d_{B(x,R)}(\mfa, \mfa_n)<2R'\}
\end{equation}
and for $j=1,2,3$
\begin{equation}
  B_j=
     \{ \mfa \in \Mf: d_{B(x,R)}(\mfa, \mfa_{m_j})<R'\}
     \, .
\end{equation}
\end{lemma}

\begin{proof}
\leanok
Let $m_1$ be the largest integer smaller than
or equal to
$n- \frac {R'} {2R}$.
Let $m_2=n$.
Let $m_3$ be the smallest integer larger than
or equal to $n+ \frac {R'} {2R}$.

Let $\mfa_{n'}\in B'$, then with \eqref{eqcarl4},
we have
\begin{equation}\label{eqcarl6}
    2R|n-n'|< 2R'\, .
\end{equation}

Assume first $n'\le m_1$. With \eqref{eqcarl6}
we have
\begin{equation*}
    R|m_1-n'|=R(m_1-n')=R(m_1-n)+R(n-n')
\end{equation*}
\begin{equation}
    < -\frac{R'}2+R'=-\frac{R'}2\, .
\end{equation}
We conclude $\mfa_{n'}\in B_1$.

Assume next $m_1<n'<m_3$. Then
$\mfa_{n'}\in B_2$.

Assume finally that $m_3\le n'$. With \eqref{eqcarl6}
we have
\begin{equation*}
    R|m_3-n'|=R(n'-m_3)=R(n'-n)+R(n-m_3)
\end{equation*}
\begin{equation}
    < R' -\frac{R'}2=-\frac{R'}2\, .
\end{equation}
We conclude $\mfa_{n'}\in B_1$.
This completes the proof of the lemma.
\end{proof}


\begin{lemma}[real van der Corput]
\label{real-van-der-Corput}
\leanok
\lean{real_van_der_Corput}
\uses{van-der-Corput}
    For any $x\in \R$ and $R>0$ and any
    function $\varphi: X\to \C$ supported on $B'=B(x,R)$
    such that
\begin{equation}
    \|\varphi\|_{\Lip(B')} = \sup_{x \in B'} |\varphi(x)| + R \sup_{x,y \in B', x \neq y} \frac{|\varphi(x) - \varphi(y)|}{\rho(x,y)}
\end{equation}
is finite and for any $n,m\in \mathbb{Z}$, we have
\begin{equation}
    \label{eq-vdc-cond1}
    \left|\int_{B'} e(\mfa_n(x)-{\mfa_m(x)}) \varphi(x) d\mu(x)\right|\le 2\pi \mu(B')\frac{\|\varphi\|_{\Lip(B')}}{1+d_{B'}(\mfa_n,\mfa_m)}
\, .
\end{equation}
\end{lemma}
\begin{proof}
    \leanok
Set $n'=n-m$. Then we have to prove
\begin{equation}
    \label{eq-vdc-cond2}
    \left|\int_{x-R}^{x+R} e^{in'y}\varphi(y) dy\right|\le 4\pi R\|\varphi\|_{\Lip(B')}
(1+2R|n'|)^{-1}\, .
\end{equation}
This follows from \Cref{van-der-Corput} with $\alpha = x - R$ and $\beta = x + R$.
\end{proof}

\begin{proof}[Proof of \Cref{real-Carleson}]
\proves{real-Carleson}
\leanok

The preceding chain of lemmas establishes that $\Mf$ is a cancellative, compatible collection of functions on $(\R, \rho, \mu, 4)$. Again, some of the statements in these lemmas are stronger than what is needed for $a=4$, but can be relaxed to give the desired conclusion for $a=4$.

With $\kappa$ as near \eqref{eq-hilker}, define
the function $K:\R\times \R\to \mathbb{C}$ as in Theorem
\ref{metric-space-Carleson} by
\begin{equation}
    K(x,y):=\kappa(x-y)\, .
\end{equation}
The function $K$ is continuous outside the diagonal
$x=y$ and vanishes on the diagonal. Hence it is measurable.


By \Cref{Hilbert-kernel-bound,Hilbert-kernel-regularity}, it follows that $K$ is a two-sided Calder\'on--Zygmund kernel on $(\R,\rho,\mu,4)$.
\Cref{Hilbert-strong-2-2} verifies \eqref{two-sided-Hr-bound-assumption}.
Thus the assumptions of \Cref{two-sided-metric-space-Carleson} are all satisfied. Applying the Theorem, \Cref{real-Carleson} follows.

\end{proof}



\begin{theorem}[classical Carleson]
    \label{classical-carleson}
    \leanok
    \lean{classical_carleson}
    \uses{exceptional-set-carleson}
    Let $f$ be a $2\pi$-periodic complex-valued continuous function on $\mathbb{R}$.
    Then for almost all $x \in \mathbb{R}$ we have
    \begin{equation}\label{eq:fourier-limit}
      \lim_{N\to\infty}S_N f(x) = f(x),
    \end{equation}
    where $S_N f$ is the $N$-th partial Fourier sum of $f$ defined in \eqref{eq:trig-series}
    with coefficients \eqref{eq:fourier-coefficients}.
\end{theorem}

\begin{proof}[Proof of \Cref{classical-carleson}]
\proves{classical-carleson}
\leanok
%Note that mere continuity implies uniform continuity in the setting of this theorem.
By applying \Cref{exceptional-set-carleson} with a sequence of $\epsilon_n:= 2^{-n}\delta$ for $n\ge 1$ and taking the union
of corresponding exceptional sets $E_n$, we see that outside a set of measure $\delta$, the partial Fourier
sums converge pointwise for $N\to \infty$. Applying this with a sequence of $\delta$ shrinking to zero and
taking the intersection of the corresponding exceptional sets, which has measure zero, we see that the Fourier
series converges outside a set of measure zero.
\end{proof}

\bibliographystyle{plain}
\bibliography{bibliography}
\end{document}
